\chapter{Về Nỗ Lực}
\markboth{Về Nỗ Lực}{On Effort}

\section{Thành công không thích người thông minh, thành công thích người kiên trì.}
\begin{truyen}{Thành công thích ai}{Who Success Likes}
\chuhoa{T}{hầy kể: ``Có hai học sinh. Một em rất nhanh trí, làm bài xong sớm rồi ngồi chơi. Em kia chậm hơn nhưng không bỏ bài khó, ngày nào cũng làm thêm vài bài. Năm sau, em chậm mà kiên trì điểm cao hơn và vững hơn. Thành công không chọn người thông minh nhất; nó chọn người không bỏ cuộc.''}
\textit{(The teacher said: ``Two students: one very quick, finished early and played. The other slower but never skipped hard problems, did extra every day. Next year the persistent one had higher and steadier grades. Success doesn't pick the smartest; it picks the one who doesn't quit.'')}
\end{truyen}

\begin{baihoc}
Kiên trì bù được thiếu thông minh; thông minh mà lười thì không đi xa.
\textit{(Persistence can make up for less smarts; smart but lazy doesn't go far.)}
\end{baihoc}

\begin{ghinhoanh}
\textbf{Short English to remember:}

Success favors the persistent, not just the smart.

\end{ghinhoanh}

\begin{tuhocnhanh}
1. Vì sao kiên trì quan trọng hơn thông minh? \textit{Why does persistence matter more than smarts?}

2. Em có đang kiên trì với việc khó không? \textit{Are you persisting with what's hard?}
\end{tuhocnhanh}
\ngancach

\section{Người giỏi là người làm lại bài sai nhiều lần nhất.}
\begin{truyen}{Làm lại bài sai}{Redoing Wrong Problems}
\chuhoa{M}{ột em hỏi: ``Thưa thầy, làm sao để giỏi như bạn A?'' Thầy đáp: ``Thầy không biết bạn A học bao nhiêu giờ. Thầy chỉ biết bạn ấy có quyển sổ chép lại mọi bài sai và tuần nào cũng làm lại. Em thử làm giống vậy — người giỏi thường là người làm lại bài sai nhiều lần nhất.''}
\textit{(A student asked: ``How can I get as good as A?'' The teacher said: ``I don't know how many hours A studies. I only know A has a notebook of every wrong problem and redoes them every week. Try that — the strong ones are usually the ones who redo wrong problems the most.'')}
\end{truyen}

\begin{baihoc}
Bài sai là tài liệu quý nhất nếu em chịu làm lại và hiểu vì sao sai.
\textit{(Wrong problems are the most valuable material if you redo them and understand why you were wrong.)}
\end{baihoc}

\begin{ghinhoanh}
\textbf{Short English to remember:}

The best learners are those who redo their mistakes the most.

\end{ghinhoanh}

\begin{tuhocnhanh}
1. Em có ghi lại và làm lại bài sai không? \textit{Do you write down and redo wrong problems?}

2. Vì sao làm lại bài sai lại giúp giỏi lên? \textit{Why does redoing mistakes make you better?}
\end{tuhocnhanh}
\ngancach

\section{Không ai leo lên đỉnh núi bằng thang máy.}
\begin{truyen}{Leo núi}{Climbing the Mountain}
\chuhoa{T}{hầy nói khi cả lớp than bài khó: ``Không ai lên đỉnh núi bằng thang máy. Ai cũng phải bước từng bước. Em thấy khó là đúng — vì em đang leo. Chỉ cần không dừng lại giữa đường; từng bước nhỏ sẽ đưa em lên.''}
\textit{(When the class complained the lesson was hard, the teacher said: ``No one reaches the top by elevator. Everyone has to take step by step. Feeling it's hard is normal — you're climbing. Just don't stop halfway; small steps will take you up.'')}
\end{truyen}

\begin{baihoc}
Đường lên đỉnh luôn là bước chân, không phải một bước nhảy; kiên trì từng bước mới tới.
\textit{(The path to the top is always footsteps, not one leap; only persistent steps get you there.)}
\end{baihoc}

\begin{ghinhoanh}
\textbf{Short English to remember:}

No one reaches the top by elevator.

\end{ghinhoanh}

\begin{tuhocnhanh}
1. Khi bài khó, em thường bỏ cuộc hay bước thêm một bước? \textit{When it's hard, do you quit or take one more step?}

2. ``Từng bước nhỏ'' trong học của em là gì? \textit{What are ``small steps'' in your learning?}
\end{tuhocnhanh}
\ngancach

\section{Mỗi lần em cố gắng thêm một chút, tương lai em đổi hướng một chút.}
\begin{truyen}{Đổi hướng tương lai}{Changing the Future's Direction}
\chuhoa{T}{hầy giải thích: ``Mỗi lần em cố thêm một chút — một bài khó không bỏ, một lần hỏi thầy, một buổi ôn thêm — tương lai em không đứng yên. Nó đổi hướng một chút. Nhiều lần đổi nhỏ cộng lại sẽ thành con đường khác hẳn.''}
\textit{(The teacher explained: ``Each time you push a little more — one hard problem you don't skip, one question you ask, one extra review — your future doesn't stay the same. It shifts a little. Many small shifts add up to a completely different path.'')}
\end{truyen}

\begin{baihoc}
Tương lai không cố định; mỗi lần cố gắng thêm là một lần nó nghiêng theo hướng tốt hơn.
\textit{(The future isn't fixed; each extra effort tilts it toward something better.)}
\end{baihoc}

\begin{ghinhoanh}
\textbf{Short English to remember:}

Every extra bit of effort shifts your future a little.

\end{ghinhoanh}

\begin{tuhocnhanh}
1. Hôm nay em đã cố thêm ``một chút'' nào chưa? \textit{Did you push ``a little more'' today?}

2. Em tin rằng cố gắng nhỏ sẽ đổi hướng tương lai không? \textit{Do you believe small efforts can change your future direction?}
\end{tuhocnhanh}
\ngancach

\section{Học sinh mạnh không phải là người không ngã, mà là người đứng dậy nhanh nhất.}
\begin{truyen}{Đứng dậy nhanh}{Getting Up Fast}
\chuhoa{S}{au khi một em bị điểm kém và ngồi buồn cả buổi, thầy nói: ``Ngã thì ai cũng ngã. Người mạnh không phải người không bao giờ ngã, mà người ngã xong đứng dậy nhanh — xem bài sai ở đâu, sửa, rồi bước tiếp. Em thử đứng dậy ngay bây giờ.''}
\textit{(After a student got a bad grade and sat sad all period, the teacher said: ``Everyone falls. The strong one isn't who never falls, but who gets up fast — sees where the mistake was, fixes it, and moves on. Try getting up right now.'')}
\end{truyen}

\begin{baihoc}
Sức mạnh không nằm ở không vấp ngã, mà ở tốc độ đứng dậy và bước tiếp.
\textit{(Strength isn't in never falling; it's in how fast you get up and keep going.)}
\end{baihoc}

\begin{ghinhoanh}
\textbf{Short English to remember:}

The strongest student is not the one who never falls, but the one who gets up fastest.

\end{ghinhoanh}

\begin{tuhocnhanh}
1. Khi em ``ngã'' (điểm kém, sai bài), em đứng dậy nhanh hay chần chừ? \textit{When you ``fall'', do you get up fast or hesitate?}

2. Làm sao để đứng dậy nhanh sau sai lầm? \textit{How can you get up quickly after a mistake?}
\end{tuhocnhanh}
\ngancach

\section{Chăm chỉ hôm nay chính là đỡ vất vả ngày mai.}
\begin{truyen}{Đỡ vất vả ngày mai}{Less Struggle Tomorrow}
\chuhoa{T}{hầy nói với những em hay để dồn bài đến gần thi: ``Chăm chỉ hôm nay không phải để thầy vui. Đó là cách em đỡ vất vả cho chính em ngày mai. Hôm nay em làm một phần, mai em bớt gánh một phần. Còn để hết đến cuối mới làm thì một đêm em gánh cả núi.''}
\textit{(The teacher told students who always left work until near exams: ``Working hard today isn't to please me. It's how you make tomorrow easier for yourself. Do a part today, and you carry less tomorrow. Leave everything to the end and you'll carry a mountain in one night.'')}
\end{truyen}

\begin{baihoc}
Chăm từng ngày là chia gánh; dồn một lúc là gánh cả núi.
\textit{(Working daily spreads the load; cramming once is carrying the whole mountain.)}
\end{baihoc}

\begin{ghinhoanh}
\textbf{Short English to remember:}

Working hard today is how you make tomorrow easier.

\end{ghinhoanh}

\begin{tuhocnhanh}
1. Em có hay dồn việc đến gần thi không? \textit{Do you often cram before exams?}

2. Làm thế nào để ``đỡ vất vả ngày mai''? \textit{How can you make tomorrow less of a struggle?}
\end{tuhocnhanh}
\ngancach

\section{Không có con đường tắt đến tri thức.}
\begin{truyen}{Không có đường tắt}{No Shortcut to Knowledge}
\chuhoa{M}{ột em hỏi: ``Thưa thầy, có cách nào học nhanh mà vẫn nhớ lâu không?'' Thầy cười: ``Nếu có thì thầy đã bán rồi. Thật ra không có đường tắt. Học đều, làm bài, sửa sai — con đường ấy dài nhưng là con đường thật. Đường tắt thường dẫn ra chỗ không có gì.''}
\textit{(A student asked: ``Is there a way to learn fast and remember long?'' The teacher smiled: ``If there were, I'd have sold it. There really is no shortcut. Study steadily, do problems, fix mistakes — that path is long but real. Shortcuts usually lead nowhere.'')}
\end{truyen}

\begin{baihoc}
Tri thức đích thực không có lối tắt; đi đúng đường dài mới tới bền.
\textit{(Real knowledge has no shortcut; the long right path is what lasts.)}
\end{baihoc}

\begin{ghinhoanh}
\textbf{Short English to remember:}

There is no shortcut to knowledge.

\end{ghinhoanh}

\begin{tuhocnhanh}
1. Em có từng tìm ``đường tắt'' trong học không? \textit{Have you ever looked for a ``shortcut'' in learning?}

2. Con đường học ``đúng'' với em là gì? \textit{What is the ``right'' learning path for you?}
\end{tuhocnhanh}
\ngancach

\section{Nếu em thấy khó, nghĩa là em đang tiến lên.}
\begin{truyen}{Thấy khó là đang tiến}{Hard Means Moving Forward}
\chuhoa{K}{hi cả lớp kêu bài khó, thầy nói: ``Nếu em thấy dễ quá, có thể em đang đứng yên. Khi em thấy khó, nghĩa là em đang bước vào vùng mới — đang tiến. Khó không phải dấu hiệu em kém; khó là dấu hiệu em đang leo.''}
\textit{(When the class said the lesson was hard, the teacher said: ``If everything feels too easy, you might be standing still. When it feels hard, you're stepping into new territory — you're advancing. Hard isn't a sign you're weak; hard is a sign you're climbing.'')}
\end{truyen}

\begin{baihoc}
Cảm giác khó thường đi kèm với việc đang học cái mới; đừng sợ khó.
\textit{(The feeling of difficulty often comes with learning something new; don't fear hard.)}
\end{baihoc}

\begin{ghinhoanh}
\textbf{Short English to remember:}

If it feels hard, you're moving forward.

\end{ghinhoanh}

\begin{tuhocnhanh}
1. Khi em thấy bài khó, em nghĩ ``mình kém'' hay ``mình đang tiến''? \textit{When work feels hard, do you think ``I'm weak'' or ``I'm moving forward''?}

2. Làm sao để không sợ cảm giác khó? \textit{How can you stop fearing the feeling of difficulty?}
\end{tuhocnhanh}
\ngancach

\section{Người thành công thường chỉ làm thêm một bước nữa so với người bỏ cuộc.}
\begin{truyen}{Một bước nữa}{One More Step}
\chuhoa{T}{hầy kể: ``Hai người cùng leo núi. Cả hai mệt, cùng muốn dừng. Một người ngồi xuống. Người kia nói: mình đi thêm một bước nữa thôi. Rồi một bước nữa. Cuối cùng người ấy lên đỉnh. Khác biệt không phải sức khỏe — chỉ là một bước nữa.''}
\textit{(The teacher said: ``Two people climb a mountain. Both are tired, both want to stop. One sits down. The other says: I'll take just one more step. Then one more. In the end that person reaches the top. The difference isn't strength — it's one more step.'')}
\end{truyen}

\begin{baihoc}
Thành công và bỏ cuộc đôi khi chỉ cách nhau một bước — bước em có chịu bước hay không.
\textit{(Success and giving up are sometimes one step apart — the step you choose to take or not.)}
\end{baihoc}

\begin{ghinhoanh}
\textbf{Short English to remember:}

Winners often do just one more step than those who quit.

\end{ghinhoanh}

\begin{tuhocnhanh}
1. Lần gần nhất em muốn bỏ cuộc, em có làm thêm ``một bước nữa'' không? \textit{Last time you wanted to quit, did you take one more step?}

2. Làm sao để nhớ ``một bước nữa'' khi mệt? \textit{How can you remember ``one more step'' when you're tired?}
\end{tuhocnhanh}
\ngancach

\section{Cố gắng hôm nay sẽ trở thành sự tự tin ngày mai.}
\begin{truyen}{Tự tin ngày mai}{Confidence Tomorrow}
\chuhoa{T}{hầy nói với em hay run khi phát biểu: ``Mỗi lần em cố gắng giơ tay, trả lời — dù sai — em đang gửi tiền vào ngân hàng tự tin. Hôm nay em cố một chút, ngày mai em tự tin thêm một chút. Cố gắng hôm nay chính là sự tự tin ngày mai.''}
\textit{(The teacher told a student who was nervous to speak: ``Every time you try to raise your hand and answer — even if wrong — you're depositing into the confidence bank. Try a little today, feel a little more confident tomorrow. Today's effort is tomorrow's confidence.'')}
\end{truyen}

\begin{baihoc}
Tự tin không tự nhiên có; nó được xây bằng từng lần cố gắng nhỏ.
\textit{(Confidence doesn't appear by itself; it's built by each small effort.)}
\end{baihoc}

\begin{ghinhoanh}
\textbf{Short English to remember:}

Today's effort becomes tomorrow's confidence.

\end{ghinhoanh}

\begin{tuhocnhanh}
1. Em có thấy mình tự tin hơn sau mỗi lần cố gắng không? \textit{Do you feel more confident after each try?}

2. Em đang ``gửi'' gì vào ``ngân hàng tự tin'' mỗi ngày? \textit{What are you ``depositing'' into your confidence bank each day?}
\end{tuhocnhanh}
\ngancach
